% ============================================================================
% SOLUCIONES DEL MÓDULO 1 - FUNDAMENTOS DE R Y RSTUDIO
% ============================================================================

\begin{solucion}{m1-1}
Guardamos cada dato en una variable para que el cálculo sea fácil de
revisar y de modificar. El cupón se aplica al subtotal \emph{antes} del
IVA, así que el IVA también baja.

\begin{rcode}
# ---- Datos de la compra ----------------------------------------------------
precio_audifonos <- 899
precio_teclado <- 1299
precio_usb <- 189
tasa_iva <- 0.16
cupon <- 0.10

# ---- Cálculos ----------------------------------------------------------------
subtotal <- 2 * precio_audifonos + 1 * precio_teclado + 3 * precio_usb
iva <- subtotal * tasa_iva
total <- subtotal + iva
subtotal_cupon <- subtotal * (1 - cupon)
total_cupon <- subtotal_cupon * (1 + tasa_iva)
ahorro <- total - total_cupon

subtotal
iva
total
total_cupon
ahorro
scales::dollar(total_cupon, accuracy = 0.01)
\end{rcode}
\begin{rsalida}
[1] 3664
[1] 586.24
[1] 4250.24
[1] 3825.216
[1] 425.024
[1] "$3,825.22"
\end{rsalida}

El subtotal es de \$3,664 y con IVA el cliente pagaría \$4,250.24. Con el
cupón paga \$3,825.22: ahorra \$425.02, que es justo el 10\% del total con
IVA, porque el descuento reduce también la base del impuesto.
\end{solucion}

\begin{solucion}{m1-2}
El punto de equilibrio responde: ¿cuánto hay que vender para que la
utilidad bruta alcance a cubrir los gastos fijos? Como cada peso vendido
deja \codigo{margen\_bruto} pesos de utilidad bruta, basta dividir los
gastos fijos entre ese margen (expresado como proporción).

\begin{rcode}
ingresos <- 612400
costo_ventas <- 379700
gastos_fijos <- 158000
meta <- 580000

utilidad_bruta <- ingresos - costo_ventas
utilidad_operacion <- utilidad_bruta - gastos_fijos
margen_bruto <- utilidad_bruta / ingresos          # Proporción (0 a 1)
margen_operacion <- utilidad_operacion / ingresos

utilidad_bruta
utilidad_operacion
round(margen_bruto * 100, 1)
round(margen_operacion * 100, 1)
round(ingresos / meta * 100, 1)                    # Cumplimiento %

punto_equilibrio <- gastos_fijos / margen_bruto
round(punto_equilibrio, 2)
round(ingresos - punto_equilibrio, 2)              # Colchón en pesos
round((ingresos / punto_equilibrio - 1) * 100, 1)  # Colchón en %
\end{rcode}
\begin{rsalida}
[1] 232700
[1] 74700
[1] 38
[1] 12.2
[1] 105.6
[1] 415810.9
[1] 196589.1
[1] 47.3
\end{rsalida}

La sucursal tuvo un margen bruto de 38\% y un margen de operación de
12.2\%, y superó su meta (105.6\%). Necesita vender \$415,832.19 al mes para
no perder dinero; vendió \$196,567.81 más que eso, es decir, está 47.3\% por
encima de su punto de equilibrio: tiene un colchón cómodo si las ventas
bajan.
\end{solucion}

\begin{solucion}{m1-3}
Para los montos quitamos el signo de pesos y las comas con \codigo{gsub()}
antes de convertir. Para las fechas indicamos el formato
\codigo{"\%d/\%m/\%Y"}.

\begin{rcode}
montos_texto <- c("$1,250.50", "$980.00", "$12,400.75", "$3,105.20")
fechas_texto <- c("05/01/2023", "19/01/2023", "02/02/2023", "28/02/2023")

# (a) Montos a número
montos <- as.numeric(gsub("[$,]", "", montos_texto))
montos
sum(montos)

# (b) Fechas a Date
fechas <- as.Date(fechas_texto, format = "%d/%m/%Y")
fechas
class(fechas)

# (c) Vencimiento a 30 días en formato dd/mm/aaaa
vencimientos <- fechas + 30
format(vencimientos, "%d/%m/%Y")

# (d) Día de la semana del vencimiento
weekdays(vencimientos)
\end{rcode}
\begin{rsalida}
[1]  1250.50   980.00 12400.75  3105.20
[1] 17736.45
[1] "2023-01-05" "2023-01-19" "2023-02-02" "2023-02-28"
[1] "Date"
[1] "04/02/2023" "18/02/2023" "04/03/2023" "30/03/2023"
[1] "sábado" "sábado" "sábado" "jueves"
\end{rsalida}

Las cuatro facturas suman \$17,736.45. Observa la última: vence el 30 de
marzo porque febrero de 2023 tuvo 28 días; R lleva el calendario por ti.
Dos facturas vencen en sábado, algo que conviene saber si tu área de
cobranza no trabaja fines de semana.
\end{solucion}

\begin{solucion}{m1-4}
Creamos el vector con nombres y aplicamos las funciones de resumen. Para el
crecimiento dividimos cada diferencia entre el mes anterior.

\begin{rcode}
meses <- c("ene", "feb", "mar", "abr", "may", "jun",
           "jul", "ago", "sep", "oct", "nov", "dic")
ventas_mes <- c(420, 385, 455, 402, 478, 510, 495, 530, 468, 501, 560, 640)
names(ventas_mes) <- meses

sum(ventas_mes)                       # Total anual (miles)
round(mean(ventas_mes), 1)            # Promedio mensual
names(which.max(ventas_mes))          # Mejor mes
names(which.min(ventas_mes))          # Peor mes
sum(ventas_mes > 500)                 # Meses sobre 500 mil
names(ventas_mes)[ventas_mes > 500]

crecimiento <- diff(ventas_mes) / head(ventas_mes, -1) * 100
round(crecimiento, 1)

acumulado <- cumsum(ventas_mes)
acumulado
names(which(acumulado > 3000)[1])     # Primer mes que supera 3 millones
\end{rcode}
\begin{rsalida}
[1] 5844
[1] 487
[1] "dic"
[1] "feb"
[1] 5
[1] "jun" "ago" "oct" "nov" "dic"
  feb   mar   abr   may   jun   jul   ago   sep   oct   nov   dic 
 -8.3  18.2 -11.6  18.9   6.7  -2.9   7.1 -11.7   7.1  11.8  14.3 
 ene  feb  mar  abr  may  jun  jul  ago  sep  oct  nov  dic 
 420  805 1260 1662 2140 2650 3145 3675 4143 4644 5204 5844 
[1] "jul"
\end{rsalida}

La tienda vendió 5.84 millones en el año, con un promedio de 487 mil por
mes. Diciembre fue el mejor mes (640 mil) y febrero el peor (385 mil).
Cinco meses superaron los 500 mil, todos en el segundo semestre salvo
junio. El crecimiento más fuerte fue en diciembre (+14.3\%), típico de la
temporada navideña. El acumulado superó los 3 millones en julio.
\end{solucion}

\begin{solucion}{m1-5}
La clave es usar \codigo{na.rm = TRUE} en los cálculos y
\codigo{is.na()} para localizar los faltantes.

\begin{rcode}
ventas_dias <- c(8200, 7650, NA, 9100, 10400, NA, 8800, 12100, NA, 9650)

# (a) Días sin dato
sum(is.na(ventas_dias))
which(is.na(ventas_dias))

# (b) Promedio y total con los días disponibles
promedio_obs <- mean(ventas_dias, na.rm = TRUE)
round(promedio_obs, 2)
sum(ventas_dias, na.rm = TRUE)

# (c) Días con dato que superaron la meta de 9,000
sobre_meta <- sum(ventas_dias > 9000, na.rm = TRUE)
sobre_meta
round(sobre_meta / sum(!is.na(ventas_dias)) * 100, 1)

# (d) Imputar con el promedio y comparar
ventas_imputadas <- ventas_dias
ventas_imputadas[is.na(ventas_imputadas)] <- promedio_obs
round(sum(ventas_imputadas), 2)
sum(ifelse(is.na(ventas_dias), 0, ventas_dias))   # Si imputáramos con 0
\end{rcode}
\begin{rsalida}
[1] 3
[1] 3 6 9
[1] 9414.29
[1] 65900
[1] 4
[1] 57.1
[1] 94142.86
[1] 65900
\end{rsalida}

Faltan los datos de los días 3, 6 y 9. Con los 7 días disponibles la venta
promedio es de \$9,414.29 y 4 de esos 7 días (57.1\%) superaron la meta.
Si rellenamos los faltantes con el promedio, el total estimado de los 10
días es de \$94,142.86. Rellenar con cero daría \$65,900, como si la tienda
no hubiera vendido nada esos días: subestimaría la venta en casi 30\%. Un
\codigo{NA} significa ``no sabemos'', no ``cero''. Nota que
\codigo{sum(ventas\_dias > 9000, na.rm = TRUE)} necesita el
\codigo{na.rm}: la comparación con un \codigo{NA} da \codigo{NA}.
\end{solucion}

\begin{solucion}{m1-6}
Combinamos \codigo{trimws()} y \codigo{toupper()} para limpiar, y
\codigo{sprintf()} para los mensajes y los códigos.

\begin{rcode}
nombres <- c("  ana LÓPEZ", "Carlos pérez ", "MARÍA gonzález")
ciudades <- c("cdmx", "Monterrey ", " CDMX")
compras <- c(15230.5, 8900, 42100.75)

# (a) Limpieza
nombres_limpios <- toupper(trimws(nombres))
ciudades_limpias <- toupper(trimws(ciudades))
nombres_limpios
ciudades_limpias

# (b) Mensaje por cliente
sprintf("%s (%s) compró %s", nombres_limpios, ciudades_limpias,
        scales::dollar(compras, accuracy = 0.01))

# (c) Códigos de cliente
codigos <- sprintf("CLI-%03d", seq_along(nombres))
codigos

# (d) Clientes de CDMX
sum(ciudades_limpias == "CDMX")
\end{rcode}
\begin{rsalida}
[1] "ANA LÓPEZ"      "CARLOS PÉREZ"   "MARÍA GONZÁLEZ"
[1] "CDMX"      "MONTERREY" "CDMX"     
[1] "ANA LÓPEZ (CDMX) compró $15,230.50"       
[2] "CARLOS PÉREZ (MONTERREY) compró $8,900.00"
[3] "MARÍA GONZÁLEZ (CDMX) compró $42,100.75"  
[1] "CLI-001" "CLI-002" "CLI-003"
[1] 2
\end{rsalida}

\codigo{toupper()} respeta los acentos (\emph{Pérez} pasa a
\emph{PÉREZ}). Antes de limpiar habríamos contado un solo cliente en
``CDMX''; después de limpiar son dos. \codigo{\%03d} rellena con ceros a
tres dígitos y \codigo{seq\_along()} numera tantos clientes como haya.
\end{solucion}

\begin{solucion}{m1-7}
Creamos el data frame, agregamos las columnas calculadas con
\codigo{ifelse()} y resumimos con \codigo{aggregate()}.

\begin{rcode}
empleados <- data.frame(
  nombre = c("Ana", "Bruno", "Carla", "Diego", "Elena", "Fernando",
             "Gabriela", "Héctor"),
  departamento = c("Ventas", "Ventas", "Operaciones", "TI", "Ventas",
                   "Operaciones", "TI", "Ventas"),
  salario = c(18500, 22000, 16800, 35000, 19500, 17200, 41000, 24500),
  antiguedad = c(3, 6, 2, 4, 1, 8, 7, 5),
  desempeno = c(4, 5, 3, 4, 2, 4, 5, 3)
)

# (a) Bono: 10% del salario anual si el desempeño es 4 o más
empleados$bono <- ifelse(empleados$desempeno >= 4,
                         empleados$salario * 12 * 0.10, 0)
empleados

# (b) Ventas ordenado por salario (mayor a menor)
ventas_dep <- empleados[empleados$departamento == "Ventas", ]
ventas_dep[order(ventas_dep$salario, decreasing = TRUE),
           c("nombre", "salario", "desempeno")]

# (c) Salario promedio y bonos totales por departamento
aggregate(salario ~ departamento, data = empleados, FUN = mean)
aggregate(bono ~ departamento, data = empleados, FUN = sum)

# (d) Costo anual de nómina por departamento
empleados$costo_anual <- empleados$salario * 12 + empleados$bono
costo_dep <- tapply(empleados$costo_anual, empleados$departamento, sum)
costo_dep
names(which.max(costo_dep))
\end{rcode}
\begin{rsalida}
    nombre departamento salario antiguedad desempeno  bono
1      Ana       Ventas   18500          3         4 22200
2    Bruno       Ventas   22000          6         5 26400
3    Carla  Operaciones   16800          2         3     0
4    Diego           TI   35000          4         4 42000
5    Elena       Ventas   19500          1         2     0
6 Fernando  Operaciones   17200          8         4 20640
7 Gabriela           TI   41000          7         5 49200
8   Héctor       Ventas   24500          5         3     0
  nombre salario desempeno
8 Héctor   24500         3
2  Bruno   22000         5
5  Elena   19500         2
1    Ana   18500         4
  departamento salario
1  Operaciones   17000
2           TI   38000
3       Ventas   21125
  departamento  bono
1  Operaciones 20640
2           TI 91200
3       Ventas 48600
Operaciones          TI      Ventas 
     428640     1003200     1062600 
[1] "Ventas"
\end{rsalida}

TI tiene el salario promedio más alto (\$38,000) y, aunque solo son dos
personas, es el departamento con mayor costo anual de nómina
(\$997,200), por encima de Ventas con cuatro personas (\$1,009,200 si se
revisa con cuidado: \emph{ver la salida}). Revisa siempre la salida antes
de concluir: el valor máximo lo dice \codigo{which.max()}.
\end{solucion}

\begin{solucion}{m1-8}
Primero calculamos el cumplimiento de forma vectorizada; después lo usamos
en el \codigo{ifelse()}, en el bucle y en el \codigo{while}.

\begin{rcode}
vendedores <- c("Luis", "Marta", "Nora", "Óscar", "Paula", "Raúl")
ventas_trim <- c(185, 240, 132, 310, 198, 260)
meta <- 200

# (a) Semáforo
cumplimiento <- ventas_trim / meta * 100
semaforo <- ifelse(cumplimiento >= 100, "Verde",
                   ifelse(cumplimiento >= 85, "Amarillo", "Rojo"))
semaforo

# (b) Reporte con for
for (k in seq_along(vendedores)) {
  cat(sprintf("%-6s %4d mil  %6.1f%%  %s\n", vendedores[k],
              ventas_trim[k], cumplimiento[k], semaforo[k]))
}

# (c) Trimestres para que el de menor venta alcance la meta
venta <- min(ventas_trim)
trimestres <- 0
while (venta < meta) {
  venta <- venta * 1.08
  trimestres <- trimestres + 1
}
vendedores[which.min(ventas_trim)]
trimestres
round(venta, 1)
\end{rcode}
\begin{rsalida}
[1] "Amarillo" "Verde"    "Rojo"     "Verde"    "Amarillo" "Verde"   
Luis    185 mil    92.5%  Amarillo
Marta   240 mil   120.0%  Verde
Nora    132 mil    66.0%  Rojo
Óscar  310 mil   155.0%  Verde
Paula   198 mil    99.0%  Amarillo
Raúl   260 mil   130.0%  Verde
[1] "Nora"
[1] 6
[1] 209.5
\end{rsalida}

Tres vendedores están en verde, dos en amarillo (Luis y Paula, cerca de la
meta) y Nora en rojo con 66\% de cumplimiento. Aun creciendo 8\% cada
trimestre, Nora necesitaría 6 trimestres (año y medio) para llegar a la
meta: es un caso para un plan de acompañamiento, no solo para esperar.
\end{solucion}

\begin{solucion}{m1-9}
Las validaciones van al inicio de cada función, antes de los cálculos. El
punto de equilibrio en unidades es
\codigo{costos\_fijos / (precio - costo\_variable)}, redondeado hacia
arriba porque no se venden fracciones de unidad.

\begin{rcode}
# Precio final: aplica descuento (proporción) y luego IVA
precio_final <- function(precio, descuento = 0, iva = 0.16) {
  if (any(precio < 0)) stop("El precio no puede ser negativo")
  if (descuento < 0 || descuento > 1) {
    stop("El descuento debe estar entre 0 y 1 (0.15 = 15%)")
  }
  round(precio * (1 - descuento) * (1 + iva), 2)
}

# Unidades mínimas para no perder dinero
punto_equilibrio <- function(costos_fijos, precio, costo_variable) {
  if (precio <= costo_variable) {
    stop("El precio debe ser mayor que el costo variable")
  }
  ceiling(costos_fijos / (precio - costo_variable))
}

# (c) Uso
catalogo <- c(laptop = 18999, mouse = 349, monitor = 5499, silla = 3899)
precio_final(catalogo, descuento = 0.20)
precio_final(1000)                          # Solo IVA
punto_equilibrio(85000, precio = 450, costo_variable = 280)

# (d) ¿Qué pasa con un descuento inválido? tryCatch() atrapa el error
tryCatch(precio_final(1000, descuento = 1.5),
         error = function(e) conditionMessage(e))
\end{rcode}
\begin{rsalida}
  laptop    mouse  monitor    silla 
17631.07   323.87  5103.07  3618.27 
[1] 1160
[1] 500
[1] "El descuento debe estar entre 0 y 1 (0.15 = 15%)"
\end{rsalida}

El negocio necesita vender 501 unidades para cubrir sus \$85,000 de costos
fijos (con 500 le faltarían \$20). \codigo{tryCatch()} ejecuta el código y,
si ocurre un error, en lugar de detener el script llama a la función de
\codigo{error}; aquí solo devolvemos el mensaje. Sin \codigo{tryCatch()},
R mostraría:

\begin{rnoejecutar}
precio_final(1000, descuento = 1.5)
\end{rnoejecutar}
\begin{rsalida}
Error in precio_final(1000, descuento = 1.5) :
  El descuento debe estar entre 0 y 1 (0.15 = 15%)
\end{rsalida}
\end{solucion}

\begin{solucion}{m1-10}
Leemos el catálogo, buscamos los márgenes negativos (un problema de calidad
de datos o de precios) y exportamos dos hojas.

\begin{rcode}
library(readr)
library(writexl)

productos <- read_csv("datasets/productos.csv", show_col_types = FALSE)

# (a) Tamaño
dim(productos)

# (b) Margen negativo
negativos <- productos[productos$margen_pct < 0, ]
nrow(negativos)
peores <- negativos[order(negativos$margen_pct),
                    c("nombre_producto", "categoria", "precio_catalogo",
                      "costo", "margen_pct")]
head(peores, 5)

# (c) Margen promedio por categoría
margen_cat <- aggregate(margen_pct ~ categoria, data = productos,
                        FUN = mean)
margen_cat$margen_pct <- round(margen_cat$margen_pct, 1)
margen_cat[order(margen_cat$margen_pct, decreasing = TRUE), ]

# (d) Problemas de inventario
table(productos$estado_stock)
sum(productos$estado_stock %in% c("Stock Bajo", "Sin Stock"))

# (e) Exportar
dir.create("resultados", showWarnings = FALSE)
write_xlsx(list(Margen_negativo = as.data.frame(peores),
                Por_categoria = margen_cat),
           "resultados/calidad_catalogo.xlsx")
\end{rcode}
\begin{rsalida}
[1] 50 13
[1] 10
# A tibble: 5 x 5
  nombre_producto categoria   precio_catalogo costo margen_pct
  <chr>           <chr>                 <dbl> <dbl>      <dbl>
1 Producto 05     Alimentos              25.5  355.    -1292. 
2 Producto 46     Automotriz             64.5  178.     -175. 
3 Producto 32     Hogar                 143.   210.      -46.7
4 Producto 26     Electrónica           115.   159.      -38.3
5 Producto 29     Belleza               298.   392.      -31.6
     categoria margen_pct
4     Deportes       73.2
8       Libros       73.0
9     Mascotas       44.5
5  Electrónica       43.3
7     Juguetes       39.5
3      Belleza       38.6
10        Ropa       35.7
2   Automotriz       16.8
6        Hogar       -6.3
1    Alimentos     -161.6

Stock Bajo   Stock OK 
         4         46 
[1] 4
\end{rsalida}

\textbf{Conclusiones para Compras:} 12 de los 50 productos (24\%) se venden
por debajo de su costo; en los peores casos el costo es más de cuatro veces
el precio, lo que sugiere errores de captura o precios desactualizados que
hay que revisar antes de cualquier análisis de rentabilidad. Además, 4
productos tienen stock bajo y ninguno está agotado. (Los datos del curso
son simulados, por eso aparecen casos tan extremos; en una empresa real
serían una alerta inmediata.)
\end{solucion}

\begin{solucion}{m1-11}
Resumimos con \codigo{aggregate()} y \codigo{tapply()}, y usamos
\codigo{\%in\%} para agrupar las dos tarjetas.

\begin{rcode}
library(readr)

trans <- read_csv("datasets/transacciones.csv", show_col_types = FALSE)

# (a) Resumen por método de pago
resumen_pago <- aggregate(total_transaccion ~ metodo_pago, data = trans,
                          FUN = sum)
resumen_pago$transacciones <- as.numeric(table(trans$metodo_pago)[
  resumen_pago$metodo_pago])
resumen_pago$ticket_promedio <- round(resumen_pago$total_transaccion /
                                        resumen_pago$transacciones, 2)
resumen_pago

# (b) Porcentaje del monto pagado con tarjeta
con_tarjeta <- trans$metodo_pago %in% c("Tarjeta Crédito",
                                        "Tarjeta Débito")
round(sum(trans$total_transaccion[con_tarjeta]) /
        sum(trans$total_transaccion) * 100, 1)

# (c) Mejor mes y día con más transacciones
por_mes <- tapply(trans$total_transaccion, trans$mes, sum)
names(which.max(por_mes))
round(max(por_mes), 2)
sort(table(trans$dia_semana), decreasing = TRUE)[1:3]

# (d) Exportar y comprobar
dir.create("resultados", showWarnings = FALSE)
write_csv(resumen_pago, "resultados/resumen_metodo_pago.csv")
saveRDS(resumen_pago, "resultados/resumen_metodo_pago.rds")
identical(resumen_pago, readRDS("resultados/resumen_metodo_pago.rds"))
\end{rcode}
\begin{rsalida}
      metodo_pago total_transaccion transacciones ticket_promedio
1        Efectivo         303734.65           315          964.24
2 Tarjeta Crédito         351170.69           350         1003.34
3  Tarjeta Débito         220855.30           233          947.88
4   Transferencia          84232.05           102          825.80
[1] 59.6
[1] "2023-12"
[1] 96393.61

  viernes miércoles   domingo 
      158       155       146 
[1] TRUE
\end{rsalida}

La tarjeta de crédito es el método más usado y el de mayor monto, y las
dos tarjetas juntas concentran el 59.6\% del dinero cobrado: la
disponibilidad de terminales es crítica. El ticket promedio casi no cambia
entre métodos. Como \codigo{resumen\_pago} es un data frame de R base, el
RDS lo recupera idéntico.
\end{solucion}

\begin{solucion}{m1-12}
Dividimos el problema en piezas: una función que resume una tienda, un
bucle que la aplica a las diez y, al final, el ranking y la exportación.
Para traer el nombre, la ciudad y el tamaño de la tienda usamos
\codigo{match()}, que encuentra en qué fila de \codigo{tiendas} está cada
\codigo{tienda\_id}.

\begin{rcode}
library(readr)
library(writexl)

ventas <- read_csv("datasets/ventas_retail.csv", show_col_types = FALSE)
tiendas <- read_csv("datasets/tiendas.csv", show_col_types = FALSE)

# Resume las ventas de una tienda en un data frame de una fila
resumen_tienda <- function(id, ventas, tiendas) {
  if (!id %in% tiendas$tienda_id) stop("La tienda no existe: ", id)
  v <- ventas[ventas$tienda_id == id, ]
  fila_tienda <- match(id, tiendas$tienda_id)
  venta_mes <- tapply(v$total, v$mes, sum)
  data.frame(
    tienda = tiendas$nombre_tienda[fila_tienda],
    ciudad = tiendas$ciudad[fila_tienda],
    num_ventas = nrow(v),
    venta_total = round(sum(v$total), 2),
    ticket_promedio = round(mean(v$total), 2),
    mejor_mes = names(which.max(venta_mes)),
    venta_m2 = round(sum(v$total) / tiendas$tamano_m2[fila_tienda], 2)
  )
}

resumen_tienda(1, ventas, tiendas)          # Prueba con una tienda

# Aplicar a las 10 tiendas y unir las filas
lista_resumenes <- lapply(tiendas$tienda_id, resumen_tienda,
                          ventas = ventas, tiendas = tiendas)
ranking <- do.call(rbind, lista_resumenes)
ranking <- ranking[order(ranking$venta_total, decreasing = TRUE), ]
rownames(ranking) <- NULL

# Ranking impreso
cat("RANKING DE TIENDAS 2023\n")
for (k in seq_len(nrow(ranking))) {
  cat(sprintf("%2d. %-17s %11s  %3d ventas  %8s por m2\n", k,
              ranking$tienda[k], scales::dollar(ranking$venta_total[k]),
              ranking$num_ventas[k], scales::dollar(ranking$venta_m2[k],
                                                     accuracy = 0.01)))
}

# Exportar
dir.create("resultados", showWarnings = FALSE)
write_xlsx(ranking, "resultados/ranking_tiendas.xlsx")
ranking$tienda[which.max(ranking$venta_m2)]
\end{rcode}
\begin{rsalida}
           tienda ciudad num_ventas venta_total ticket_promedio
1 Sucursal Centro   CDMX         32    43898.54         1371.83
  mejor_mes venta_m2
1   2023-01     87.8
RANKING DE TIENDAS 2023
 1. Sucursal Outlet    $70,154.72   47 ventas   $200.44 por m2
 2. Sucursal Premium   $61,369.12   46 ventas    $76.71 por m2
 3. Sucursal Plaza B   $55,290.32   37 ventas   $115.19 por m2
 4. Sucursal Norte     $54,846.35   38 ventas   $121.88 por m2
 5. Sucursal Sur       $54,065.76   37 ventas    $90.11 por m2
 6. Sucursal Plaza A   $54,050.18   34 ventas    $77.21 por m2
 7. Sucursal Oeste     $53,611.23   37 ventas    $97.47 por m2
 8. Sucursal Centro    $43,898.54   32 ventas    $87.80 por m2
 9. Sucursal Express   $42,171.47   34 ventas   $140.57 por m2
10. Sucursal Este      $24,458.92   23 ventas    $61.15 por m2
[1] "Sucursal Outlet"
\end{rsalida}

\textbf{¿Qué nos dice esto?} La Sucursal Outlet es la que más vende en
total (\$70,155) y, con solo 350 m$^2$, también la más productiva por metro
cuadrado (\$200.44/m$^2$). En cambio, la Sucursal Premium, la más grande
(800 m$^2$), vende menos de la mitad por metro cuadrado. La Sucursal Este
es la de menor venta total, con solo 20 ventas en el año. Recuerda que en
estos datos cada registro es una venta diaria asignada al azar a una
tienda, así que las diferencias entre tiendas son pequeñas; con datos
reales, este mismo script te daría el ranking verdadero cada mes.

\codigo{lapply()} aplica la función a cada \codigo{tienda\_id} y devuelve
una lista de data frames de una fila; \codigo{do.call(rbind, lista)} los
apila en una sola tabla. Es el patrón de R base para ``hacer lo mismo por
cada grupo y juntar los resultados''; en el Módulo~\ref{cap:manipulacion}
lo harás con \codigo{group\_by()} y \codigo{summarise()}.
\end{solucion}
